\begin{prob}
    \textbf{Follow-up: Bellman-Ford example from Lecture 14 slide 46.}

    Consider the directed weighted graph from Lecture 14 (slide 46) with source node $v_1$.
    The weighted edges are:
    \[
        (v_1,v_2):2,\;
        (v_2,v_3):-1,\;
        (v_3,v_4):1,\;
        (v_4,v_5):3,\;
        (v_5,v_6):2,\;
        (v_1,v_7):7,\;
        (v_7,v_6):0,\;
        (v_7,v_5):3,\;
        (v_5,v_7):-1.
    \]

    \textbf{Edge ordering convention.}
    In each iteration (pass), perform \texttt{update(u, v)} exactly as listed below.
    Iteration pass uses this exact edge order:
    \[
        (v_3,v_4),\; (v_1,v_2),\; (v_2,v_3),\; (v_7,v_6),\; (v_5,v_7),\;
        (v_7,v_5),\; (v_4,v_5),\; (v_5,v_6),\; (v_1,v_7).
    \]

    Answer the following:
    \begin{subprobset}
        \begin{subprob}
            How many full passes are needed until distances no longer change?
        \end{subprob}
        \begin{subprob}
            What are the final \texttt{est} and \texttt{predecessor} dictionaries?
        \end{subprob}
        \begin{subprob}
            Briefly explain why this ordering is slower.
        \end{subprob}
    \end{subprobset}

    \begin{responsebox}{4.2in}
    \end{responsebox}

    \begin{soln}
        Pass-by-pass updates under the required order:
        \begin{itemize}
            \item Pass 1: $v_2 \leftarrow 2$, $v_3 \leftarrow 1$, $v_7 \leftarrow 7$.
            \item Pass 2: $v_4 \leftarrow 2$, $v_6 \leftarrow 7$, $v_5 \leftarrow 5$.
            \item Pass 3: $v_7 \leftarrow 4$ (via $v_5 \to v_7$).
            \item Pass 4: $v_6 \leftarrow 4$ (via updated $v_7 \to v_6$).
            \item Pass 5: no values change (converged).
        \end{itemize}

        So the first no-change pass is pass 5 (equivalently, all final values are already fixed by the end of pass 4).

        Final values:
        \begin{itemize}
            \item est: $v_1\!:\!0,\; v_2\!:\!2,\; v_3\!:\!1,\; v_4\!:\!2,\; v_5\!:\!5,\; v_6\!:\!4,\; v_7\!:\!4$
            \item predecessor: $v_1\!:\!\texttt{None},\; v_2\!:\!v_1,\; v_3\!:\!v_2,\; v_4\!:\!v_3,\; v_5\!:\!v_4,\; v_6\!:\!v_7,\; v_7\!:\!v_5$
        \end{itemize}

        Why slower: this edge ordering is not aligned with the long shortest-path chain from $v_1$ to $v_6$.
        Improvements found later in a pass often cannot propagate until the next pass, so information moves more slowly.
        If we used an order that follows the chain direction (topologically along that path), the same information could propagate in far fewer passes (often around 1--2 passes for this type of example).
    \end{soln}
\end{prob}
